\section{Introdução}
\label{sec:introducao}
\aescrever
% Modelo PPGCONT: apresentar o assunto, o tema e a teoria de base.
% O objetivo NÃO aparece aqui; ele tem subseção própria.
% Fonte de conteúdo: docs/proposta/proposta-dissertacao-v5-2026-11-14.md, seção 1.1.


\subsection{Problema de pesquisa}
\label{sec:problema}
\aescrever
% Modelo: contextualizar de forma robusta e fechar com a questão-problema.
% Fonte: proposta v5, seção 1.2.


\subsection{Objetivos}
\label{sec:objetivos}
\aescrever
% Modelo: uma frase dizendo o que a seção apresenta.


\subsubsection{Objetivo Geral}
\aescrever
% Modelo: breve introdução e o objetivo geral. Fonte: proposta v5, seção 2.1.


\subsubsection{Objetivos Específicos}
\aescrever
% Modelo: no mínimo três, verbo no infinitivo e em minúscula, ponto e vírgula
% ao final de cada item e ponto final no último. Fonte: proposta v5, seção 2.2.
%
% \begin{enumerate}
%   \item verbo ...;
%   \item verbo ...;
%   \item verbo ... .
% \end{enumerate}


\subsection{Justificativa do estudo}
\label{sec:justificativa}
\aescrever
% Modelo: justificar a escolha do tema, com referências, cobrindo:
% a) originalidade; b) oportunidade; c) contributividade; d) relevância.
% Fonte: proposta v5, seção 1.3.


\subsection{Delimitação do estudo}
\label{sec:delimitacao}
\aescrever
% Fonte: proposta v5, seção 1.4.
